\documentclass{article}
\usepackage{lpsb}
\usepackage{amsmath}
\usepackage{graphicx}
\usepackage{hyperref}

\title{LPSB v3.0 Test Document}
\author{Test Author}
\date{\today}

\begin{document}
\maketitle

\begin{abstract}
This document tests all PDF 2.0 structure elements supported by LPSB v3.0.
It includes paragraphs, headings, lists, tables, figures, equations, and links.
\end{abstract}

\section{Introduction}

This is the first paragraph of the introduction section.
It should be tagged as a P element within the section.

This is the second paragraph.
LPSB should detect the paragraph boundary via the par hook.

\subsection{Background}

A subsection creates an H3 heading in PDF 2.0 terms.
This paragraph belongs to the subsection context.

\subsubsection{Details}

An H4 heading from subsubsection.

\section{Lists}

\subsection{Itemize List}

\begin{itemize}
\item First item in unordered list
\item Second item with more text
\item Third item
\end{itemize}

\subsection{Enumerate List}

\begin{enumerate}
\item First numbered item
\item Second numbered item
\item Third numbered item
\end{enumerate}

\subsection{Description List}

\begin{description}
\item[Term A] Definition of term A
\item[Term B] Definition of term B
\end{description}

\section{Tables}

\begin{table}[h]
\centering
\begin{tabular}{|l|c|r|}
\hline
Left & Center & Right \\
\hline
A1 & B1 & C1 \\
A2 & B2 & C2 \\
\hline
\end{tabular}
\caption{A simple table with three columns.}
\label{tab:simple}
\end{table}

Table~\ref{tab:simple} shows a basic table structure.

\section{Figures}

\begin{figure}[h]
\centering
% Placeholder for image - uses rule as stand-in
\rule{4cm}{3cm}
\caption{A placeholder figure.}
\label{fig:placeholder}
\end{figure}

See Figure~\ref{fig:placeholder} for the visual element.

\section{Mathematics}

Inline math: $E = mc^2$ is Einstein's famous equation.

Display equation:
\begin{equation}
\int_0^\infty e^{-x^2} dx = \frac{\sqrt{\pi}}{2}
\label{eq:gaussian}
\end{equation}

Aligned equations:
\begin{align}
a &= b + c \\
d &= e \cdot f
\end{align}

\section{Text Formatting}

This paragraph contains \textbf{bold text} and \textit{italic text}.
It also has \emph{emphasized text} for semantic markup.

\section{Links and References}

Visit \href{https://example.com}{Example Website} for more information.
This references Equation~\ref{eq:gaussian}.

\section*{Unnumbered Section}

This is an unnumbered section (starred).
It should still be detected as a heading.

\section{Conclusion}

This concludes the test document.
All structure elements should be captured in the JSON output.

\end{document}
